\subsection{Principle of localization. Locally negligible sets}

PROPOSITION 4 (Principle of localization). --- Let f be a mapping of X into a topological space F. Suppose that for every \(x \in X\) , there exist an integrable neighborhood \(V_{x}\) of x and a measurable mapping \(g_{x}\) of X into F such that \(f(y) = g_{x}(y)\) almost everywhere in \(V_{x}\) . Then f is measurable.

Let K be a compact set; there exists a finite number of points \(x_{i} \in K\) such that the \(V_{x_{i}}\) form a covering of K. It follows at once (§4, No.~9, Lemma) that there exist a negligible set \(N \subset K\) and a finite partition of K - N formed of integrable sets \(M_{j}\) such that each of the sets \(K \cap V_{x_{i}}\) is the union of a subset of N and a certain number of the \(M_{j}\) , and such that on each of the \(M_{j}\) , f is equal to one of the functions \(g_{x_{i}}\) . Now, for each \(M_{j}\) there exist a negligible set \(N_{j} \subset M_{j}\) and a partition of \(M_{j} - N_{j}\) formed by a sequence of compact sets \(K_{nj}\) ( \(n \in N\) ); on the other hand, for each \(K_{nj}\) there exist a negligible set \(P_{nj} \subset K_{nj}\) and a partition of \(K_{nj} - P_{nj}\) formed by a sequence of compact sets \(K_{mnj}\) ( \(m \in N\) ) such that the restriction of f to each of the \(K_{mnj}\) is continuous. Since the union of N, the \(N_{j}\) and the \(P_{nj}\) is negligible, f is measurable.

The concept of measurable function is therefore a concept of local character.

DEFINITION 3. --- A set \(A \subset X\) is said to be locally negligible (for the measure \(\mu\) ) if, for every \(x \in X\) , there exists a neighborhood V of x such that \(V \cap A\) is negligible.

By the principle of localization, every locally negligible set is measurable. The properties of negligible sets ( \(\S2\) ) show that every subset of a locally negligible set is locally negligible, and that every countable union of locally negligible sets is locally negligible.

PROPOSITION 5. --- For a set A to be locally negligible, it is necessary and sufficient that, for every compact set K, A∩K be negligible.

The condition is obviously sufficient since every point of X has a compact neighborhood. It is necessary, because if, for every \(x \in K\) , there exists a neighborhood \(V_{x}\) of x such that \(A \cap V_{x}\) is negligible, then there exists a finite number of points \(x_{i} \in K\) such that the \(V_{x_{i}}\) form a covering of K, and \(A \cap K\) is contained in the union of the negligible sets \(A \cap V_{x_{i}}\) .

COROLLARY 1. --- For a set A to be negligible, it is necessary and sufficient that it be locally negligible and of finite outer measure.

The condition is obviously necessary. Conversely, if it is satisfied then A is contained in an integrable open set G, which is the union of a negligible set N and a sequence \((\mathrm{K}_{n})\) of compact sets ( \(\S4\) , No.~6, Cor. 2 of Th. 4); since \(A \cap N\) and the sets \(A \cap K_{n}\) are negligible, the same is true of their union A.

COROLLARY 2. --- Every locally negligible open set is negligible (and is therefore contained in the complement of the support of \(\mu\) ).

For, the outer measure of an open set G is the supremum of the measures \(|\mu|(\mathrm{K})\) of the compact sets \(K \subset G\) ( \(\S4\) , No.~6, Cor. 4 of Th. 4); if G is locally negligible then \(|\mu|(\mathrm{K}) = 0\) for every compact set K contained in G, therefore \(|\mu|^{*}(\mathrm{G}) = 0\) .

COROLLARY 3. --- In a locally compact space X that is countable at infinity, every locally negligible set is negligible.

Since X is the union of a sequence \((\mathrm{K}_{n})\) of compact sets, every locally negligible set A is the union of the negligible sets \(A\cap K_{n}\) , hence is negligible.

One can give examples of locally compact spaces that are not countable at infinity, and of measures on such a space X such that there exist sets in X that are locally negligible but not negligible (§1, Exer. 5).

COROLLARY 4. --- Let f be a mapping of X into a topological space F. If the set N of points of discontinuity of f is locally negligible, then f is measurable.

For every compact set \(K \subset X\) , \(K \cap N\) is negligible (Prop. 5), therefore, for every \(\varepsilon > 0\) , there exists a compact set \(K_{1} \subset K - (K \cap N)\) such that \(|\mu|(\mathrm{K} - \mathrm{K}_{1}) \leqslant \varepsilon\) (§4, No.~6, Th. 4), and by hypothesis the restriction of f to \(K_{1}\) is continuous, whence the conclusion.

If \(P\{x\}\) is a property, the property \(\langle P\{x\}\) locally almost everywhere (with respect to \(\mu\) ) is by definition equivalent to the property \(\langle\) the set of x such that \((x \in X\) and not \(P\{x\}\) ) is locally negligible (for the measure \(\mu\) ) . If F is any set, the relation \(\langle f(x) = g(x)\) locally almost everywhere is an equivalence relation in the set of mappings of X into F. In particular, if F is a vector space, a mapping f of X into F such that \(f(x) = 0\) locally almost everywhere is said to be locally negligible. We leave to the reader the task of establishing for these concepts most of the properties corresponding to those that have been enumerated in §2, Nos. 4, 5 and 6 for functions equal almost everywhere. We shall limit ourselves to observing that if two continuous mappings \(f, g\) of X into a Hausdorff topological space F are equal locally almost everywhere, then they are equal almost everywhere by virtue of Cor. 2 of Prop. 5 (hence are equal at every point of the support of \(\mu\) (§2, No.~4, Prop. 9)); however, we state explicitly the following proposition, which is an immediate consequence of the principle of localization:

PROPOSITION 6. --- Let f be a measurable mapping of X into a topological space F. Every mapping of X into F, equal to f locally almost everywhere, is measurable.

